\part{L'Interfaccia}

\chapter{Schermata Principale}

FairWindSK mantiene sempre visibili tre zone della shell, indipendentemente dall'applicazione
attiva. Questi nomi sono usati in modo uniforme in questo manuale, nel codice sorgente del
progetto e in tutta la documentazione per i collaboratori.

\section{Le Tre Zone}

\begin{center}
\begin{tcolorbox}[enhanced,colback=LtBlue,colframe=Navy,boxrule=1.5pt,arc=4pt,
  top=7pt,bottom=5pt,left=12pt,right=12pt,width=0.96\textwidth]
  {\sffamily\bfseries\color{Navy}BARRA SUPERIORE --- sempre visibile}\\[3pt]
  {\small\icon{lcd-app}~Icona app \quad
         \icon{layout-signalk}~Widget dati \quad
         Nome app \quad
         \icon{lcd-clock}~Data/ora \quad
         Icone di stato}
\end{tcolorbox}
\vspace{2pt}
\begin{tcolorbox}[enhanced,colback=white,colframe=Navy,boxrule=1.5pt,arc=4pt,
  top=26pt,bottom=26pt,left=12pt,right=12pt,width=0.96\textwidth]
  \centering
  {\large\sffamily\bfseries\color{MidGray}AREA APPLICAZIONI}\\[4pt]
  {\small\sffamily\color{MidGray}\itshape
    Launcher \textbar{} Plotter \textbar{} Strumenti \textbar{}
    Impostazioni \textbar{} MyData \textbar{} Applicazioni web Signal~K}
\end{tcolorbox}
\vspace{2pt}
\begin{tcolorbox}[enhanced,colback=LtBlue,colframe=Navy,boxrule=1.5pt,arc=4pt,
  top=7pt,bottom=5pt,left=12pt,right=12pt,width=0.96\textwidth]
  {\sffamily\bfseries\color{Navy}BARRA INFERIORE --- sempre visibile}\\[3pt]
  {\small
    Striscia app aperte \quad
    \icon{applications}~App \quad
    \icon{alarm-pob}~POB \quad
    \icon{command-autopilot}~Autopilota \quad
    \icon{anchor-iec}~Ancora \quad
    \icon{alarm-general}~Allarmi \quad
    \icon{database}~MyData \quad
    \icon{settings-iec}~Impostazioni \quad
    \icon{layout-signalk}~Stato}
\end{tcolorbox}
\end{center}

\screenshot{screen-layout-overview}{Schema schermata: Barra Superiore, Area Applicazioni, Barra Inferiore}{}

\section{L'Area Applicazioni}

L'Area Applicazioni è lo spazio di lavoro principale tra le barre. Qui FairWindSK mostra:

\begin{itemize}
  \item le pagine del launcher
  \item app di carteggio e navigazione
  \item Impostazioni
  \item MyData
  \item applicazioni web Signal~K incorporate
\end{itemize}

\section{Cassetti (Drawer)}

\rowcolors{2}{LtGray}{white}
\begin{longtable}{L{4.0cm}L{9.7cm}}
\toprule
\rowcolor{Navy}
\color{white}\sffamily\bfseries Cassetto &
\color{white}\sffamily\bfseries Descrizione \\
\midrule
Cassetto orizzontale Barra Inferiore &
  Dialogo modale che si apre verso l'alto dalla Barra Inferiore.
  Usato per dialoghi modali web, conferme, selettori file e pannelli di impostazioni.
  Un solo dialogo alla volta; dialoghi aggiuntivi vengono accodati.
  Può estendersi verso l'alto fino a coprire temporaneamente l'Area Applicazioni. \\
Cassetto verticale Area Applicazioni &
  Pannello non modale che si apre verso sinistra dal bordo destro dell'Area Applicazioni.
  Usato per menu contestuali o palette di strumenti.
  L'Area Applicazioni rimane interattiva dietro di esso. \\
\bottomrule
\end{longtable}

\section{Modello a Finestra Singola}

FairWindSK applica un modello a finestra singola su tutte le piattaforme.

Android costituisce l'eccezione intenzionale per le applicazioni native selezionate.
Da Android 13/API 33, FairWindSK può essere scelto facoltativamente come app Home e
può avviare esplicitamente le attività selezionate in \ui{Impostazioni~> Android}.
Le attività native lasciano temporaneamente la finestra FairWindSK e vi ritornano
tramite Indietro o Home; i contenuti FairWindSK e Signal~K restano nel modello a
finestra singola.
Nessun contenuto FairWindSK apre una seconda finestra; fuori dall'eccezione Android
descritta sopra, non viene avviata alcuna applicazione nativa esterna.
Le voci del launcher con schema \code{file://} sono bloccate.

Su piattaforme desktop (Linux, macOS, Windows, Raspberry~Pi~OS), premi \key{SHIFT+TAB}
per riportare FairWindSK in primo piano se un'applicazione web ospitata ha acquisito
il focus completo.

\chapter{La Barra Superiore}
\label{ch:barrasuperiore}

La Barra Superiore è sempre visibile.
Il suo scopo: fornire un controllo di fiducia in due secondi in qualsiasi momento.
Usala per rispondere a: \textit{Cosa sta succedendo adesso, e i dati sono ancora affidabili?}

\screenshot{topbar-overview}{Barra Superiore: posizione, SOG, DPT, BTW, ETA e icone di stato}{}

\section{Elementi Fissi della Barra Superiore}

\rowcolors{2}{LtGray}{white}
\begin{longtable}{C{1.5cm}L{2.8cm}L{9.4cm}}
\toprule
\rowcolor{Navy}
\color{white}\sffamily\bfseries Icona &
\color{white}\sffamily\bfseries Elemento &
\color{white}\sffamily\bfseries Funzione \\
\midrule
\iconlg{lcd-app}       & Icona app    &
  Icona dell'applicazione corrente.
  Toccarla apre il cassetto verticale dell'Area Applicazioni. \\
---                     & Nome app     &
  Nome dell'applicazione attualmente attiva. \\
\iconlg{lcd-clock}     & Data/ora     &
  Data e ora corrente del dispositivo. \\
\iconlg{layout-signalk}& Icone stato  &
  Indicatori di salute REST e stream. \\
\bottomrule
\end{longtable}

\section{Widget Dati --- Riferimento Completo}

Ogni widget si abbona a un percorso Signal~K configurato in
\ui{Impostazioni~> Percorsi Signal~K}.

\rowcolors{2}{LtGray}{white}
\begin{longtable}{C{1.4cm}C{1.1cm}L{4.4cm}L{3.0cm}L{3.8cm}}
\toprule
\rowcolor{Navy}
\color{white}\sffamily\bfseries Icona &
\color{white}\sffamily\bfseries Chiave &
\color{white}\sffamily\bfseries Percorso Signal~K predefinito &
\color{white}\sffamily\bfseries Visualizza &
\color{white}\sffamily\bfseries Rilevanza operativa \\
\midrule
\iconlg{lcd-position} & \code{pos} & \skpath{navigation.position}                              & Lat / Lon             & Conferma fix GPS \\
\iconlg{lcd-cog}      & \code{cog} & \skpath{navigation.courseOverGroundTrue}                  & COG (\textdegree{}V)  & Rotta reale sul fondale \\
\iconlg{lcd-sog}      & \code{sog} & \skpath{navigation.speedOverGround}                       & SOG (nd)              & Velocità reale per ETA \\
\iconlg{lcd-hdg}      & \code{hdg} & \skpath{navigation.headingTrue}                           & Prua (\textdegree{}V) & Dove punta la prua \\
\iconlg{lcd-stw}      & \code{stw} & \skpath{navigation.speedThroughWater}                     & STW (nd)              & Assetto vele, deduzione corrente \\
\iconlg{lcd-dpt}      & \code{dpt} & \skpath{environment.depth.belowTransducer}                & Profondità            & Critico in acque basse \\
\iconlg{lcd-btw}      & \code{btw} & \skpath{navigation.course.calcValues.bearingTrue}         & Rilevamento WPT       & Rotta da seguire \\
\iconlg{lcd-dtg}      & \code{dtg} & \skpath{navigation.course.calcValues.distance}            & Distanza da percorrere & Progresso traversata \\
\iconlg{lcd-ttg}      & \code{ttg} & \skpath{navigation.course.calcValues.timeToGo}            & Tempo rimanente       & Pianificazione ETA \\
\iconlg{lcd-eta}      & \code{eta} & \skpath{navigation.course.calcValues.estimatedTimeOfArrival} & Ora stim.\ di arrivo & Notifica equipaggio \\
\iconlg{lcd-xte}      & \code{xte} & \skpath{navigation.course.calcValues.crossTrackError}     & Errore di rotta       & Rispetto del percorso \\
\iconlg{lcd-vmg}      & \code{vmg} & \skpath{performance.velocityMadeGood}                     & Velocità eff.\ progresso & Efficienza reale di vela \\
\iconlg{lcd-wpt}      & \code{wpt} & \skpath{navigation.course.nextPoint}                      & Nome waypoint attivo  & Consapevolezza di rotta \\
\bottomrule
\end{longtable}

\begin{warningbox}
Profondità mancante \textbf{non} significa profondità zero.
Un valore mancante significa che la misura non è disponibile --- trattala come
completamente sconosciuta.
Se la profondità diventa obsoleta o mancante in acque basse o ristrette, verifica
immediatamente e adegua di conseguenza velocità e margine di manovra.
\end{warningbox}

\section{Configurazione della Barra Superiore}

Apri \ui{Impostazioni~> Barra Superiore}.
In alto appare un'anteprima in tempo reale della barra; sotto si trova la palette dei widget.
Usa i pulsanti \iconlg{arrow-left-google}~sposta-sinistra e \iconlg{arrow-right-google}~sposta-destra
per riposizionare i widget.
Attiva/disattiva \key{Mostra Icona}, \key{Mostra Testo}, \key{Mostra Unità} e \key{Mostra Tendenza}
per ogni widget.
Tocca \iconlg{refresh-google}~\key{Ripristina Predefiniti} per ripristinare il layout predefinito.

\begin{tipbox}
In condizioni difficili, una Barra Superiore con 3--4 widget chiave (posizione, SOG, DPT, BTW)
è più sicura di dodici elementi.
Aggiungi widget solo se aggiungono consapevolezza situazionale reale nella fase di navigazione corrente.
\end{tipbox}

\chapter{La Barra Inferiore}
\label{ch:barrainferiore}

La Barra Inferiore è la striscia di azioni permanente in fondo a ogni schermata.
La sua stabilità spaziale è intenzionale: il pulsante \iconlg{applications}~App è sempre
nello stesso posto, il pulsante \iconlg{alarm-pob}~POB è sempre a portata di mano.
La memoria muscolare è sicurezza.

\screenshot{bottombar-overview}{Barra Inferiore: striscia app aperte, icone principali, stato Signal~K}{}

\section{Aree della Barra Inferiore}

\rowcolors{2}{LtGray}{white}
\begin{longtable}{L{3.6cm}L{10.1cm}}
\toprule
\rowcolor{Navy}
\color{white}\sffamily\bfseries Area &
\color{white}\sffamily\bfseries Descrizione \\
\midrule
Striscia app web aperte &
  Riga scorrevole di piccole icone app a sinistra, una per applicazione caricata.
  Tocca un'icona per passare immediatamente a quella app.
  L'applicazione attiva è evidenziata. \\
Icone principali &
  Il gruppo centrale di sette icone fisse. Stabile spazialmente in tutte le schermate. \\
Area stato Signal~K &
  Sul lato destro. Indicatori compatti di salute stream e REST, nome server,
  timestamp ultimo aggiornamento e stato della connessione.
  Toccare apre un cassetto di dettaglio. \\
\bottomrule
\end{longtable}

\section{Riferimento Icone Principali}

Queste sono le icone esatte definite in \code{BottomBar.ui}.

\rowcolors{2}{LtGray}{white}
\begin{longtable}{C{1.4cm}L{2.4cm}L{5.4cm}L{4.5cm}}
\toprule
\rowcolor{Navy}
\color{white}\sffamily\bfseries Icona &
\color{white}\sffamily\bfseries Etichetta &
\color{white}\sffamily\bfseries Funzione &
\color{white}\sffamily\bfseries Note \\
\midrule
\iconlg{applications}     & App         &
  Apre il launcher. Nasconde l'app corrente e mostra la griglia tile. &
  Uso come via di fuga sicura da qualsiasi schermata bloccata o confusa. \\
\iconlg{alarm-pob}        & POB         &
  Persona in mare. Un tocco avvia l'emergenza. &
  Sempre visibile. Mai disabilitato. \\
\iconlg{command-autopilot}& Autopilota  &
  Apre la barra Autopilota. &
  Grigio/inattivo se non viene rilevato nessun plugin compatibile. \\
\iconlg{anchor-iec}       & Ancora      &
  Apre la barra Guardia Ancora. &
  Grigio/inattivo se non viene rilevato nessun plugin ancora. \\
\iconlg{alarm-general}    & Allarmi     &
  Apre la barra Allarmi. Evidenziato quando un allarme è attivo. &
  Colore accent quando attivo; attira l'attenzione su tutti i preset. \\
\iconlg{database}         & MyData      &
  Apre il gestore risorse (Cap.~\ref{ch:mydata}). & \\
\iconlg{settings-iec}     & Impostazioni &
  Apre tutte le configurazioni (Cap.~\ref{ch:impostazioni}). & \\
\bottomrule
\end{longtable}

\section{Barre Operative Sovrapposte}

Toccando \iconlg{alarm-pob}~POB, \iconlg{command-autopilot}~Autopilota,
\iconlg{anchor-iec}~Ancora, o \iconlg{alarm-general}~Allarmi si sostituisce la Barra
Inferiore normale con una barra operativa dedicata.
L'Area Applicazioni continua a funzionare in background.
Tocca \iconlg{close-google}~chiudi (icona più a destra) per tornare alla Barra Inferiore
normale senza modificare lo stato operativo.

\begin{warningbox}
Le barre operative non sono modali.
I dati continuano ad aggiornarsi.
Il plotter e tutte le applicazioni rimangono attivi in background durante l'uso delle barre operative.
\end{warningbox}

\chapter{Il Launcher e la Gestione delle Applicazioni}
\label{ch:launcher}

Il launcher è la posizione di partenza sicura della shell.
Usalo per tornare alle tile app, passare da un'attività all'altra, uscire da un'app bloccata
o riorientarti dopo un flusso confuso.
Quando la schermata corrente smette di sembrare chiara o affidabile, tocca
\iconlg{applications}~\key{App} per tornare qui.

\screenshot{launcher-overview}{Schermata principale del launcher: griglia tile, indicatore pagina, Barra Inferiore}{}

\section{Aprire un'Applicazione}

\begin{enumerate}
  \item Tocca \iconlabel{applications}{App} nella Barra Inferiore.
  \item Se necessario, scorri tra le pagine.
  \item Tocca la tile dell'applicazione desiderata.
  \item Attendi che si carichi nell'Area Applicazioni.
\end{enumerate}

Se un'app non si presenta correttamente, evita tocchi ripetuti rapidi.
Attendi brevemente, riprova una volta se la shell lo propone, poi torna al launcher
se l'app sembra ancora anomala.

\section{Passare tra Applicazioni Aperte}

Usa la striscia app aperte nella Barra Inferiore per passare tra le applicazioni caricate
con un singolo tocco.

\section{Aggiungere un'Applicazione Personalizzata}

\begin{enumerate}
  \item Apri \ui{Impostazioni~> Applicazioni}.
  \item Tocca \iconlabel{define-new-application-in-palette}{Crea app}.
  \item Inserisci un nome visualizzato e l'URL completo (\code{http://} o \code{https://}).
  \item Imposta facoltativamente un percorso icona personalizzato e il livello di zoom.
  \item Tocca \key{Salva}.
\end{enumerate}

\section{Ripristinare un'Applicazione Bloccata o Vuota}

\begin{enumerate}
  \item Attendi brevemente. Riprova una volta se la shell offre un'azione di ripresa.
  \item Tocca \iconlabel{applications}{App} per tornare al launcher.
  \item Controlla la salute della shell e lo stato della connessione.
  \item Riapri l'app solo dopo aver confermato che la shell stessa è integra.
\end{enumerate}

\begin{tipbox}
Un'applicazione in errore non significa che FairWindSK abbia smesso di funzionare.
La shell, i widget dati e le barre operative (POB, ancora, autopilota) continuano
a funzionare indipendentemente da ogni singola applicazione.
Un'app degradata non è la stessa cosa di uno stato dell'imbarcazione degradato.
\end{tipbox}
